\section{Mittag-Leffler modules}
\label{section-mittag-leffler-modules}

\noindent
A Mittag-Leffler module is (very roughly) a module which can be written
as a directed limit whose dual is a Mittag-Leffler system. To be able
to give a precise definition we need to do a bit of work.

\begin{definition}
\label{definition-ML-inductive-system}
Let $(M_i, f_{ij})$ be a directed system of $R$-modules.  We say that
$(M_i, f_{ij})$ is a {\it Mittag-Leffler directed system of modules} if each
$M_i$ is an $R$-module of finite presentation and if for every $R$-module $N$,
the inverse system
$$
(\Hom_R(M_i, N), \Hom_R(f_{ij}, N))
$$
is Mittag-Leffler.
\end{definition}

\noindent
We are going to characterize those $R$-modules that are colimits of
Mittag-Leffler directed systems of modules.

\begin{definition}
\label{definition-domination}
Let $f: M \to N$ and $g: M \to M'$ be maps of $R$-modules.
Then we say $g$ {\it dominates} $f$ if for any $R$-module $Q$, we have $\Ker(f
\otimes_R \text{id}_Q) \subset \Ker(g \otimes_R \text{id}_Q)$.
\end{definition}

\noindent
It is enough to check this condition for finitely presented modules.

\begin{lemma}
\label{lemma-domination-fp}
Let $f: M \to N$ and $g: M \to M'$ be maps of $R$-modules.
Then $g$ dominates $f$ if and only if for any finitely presented $R$-module
$Q$, we have $\Ker(f \otimes_R \text{id}_Q) \subset \Ker(g \otimes_R
\text{id}_Q)$.
\end{lemma}

\begin{proof}
Suppose $\Ker(f \otimes_R \text{id}_Q) \subset \Ker(g \otimes_R
\text{id}_Q)$ for all finitely presented modules $Q$.  If $Q$ is an
arbitrary module, write $Q = \colim_{i \in I} Q_i$ as a colimit of a
directed
system of finitely presented modules $Q_i$.  Then $\Ker(f \otimes_R
\text{id}_{Q_i}) \subset \Ker(g \otimes_R \text{id}_{Q_i})$ for
all $i$.  Since taking directed colimits is exact and commutes with tensor
product, it follows that $\Ker(f \otimes_R \text{id}_Q) \subset \Ker(g
\otimes_R \text{id}_Q)$.
\end{proof}

\begin{lemma}
\label{lemma-domination-universally-injective}
Let $f : M \to N$ and $g : M \to M'$ be maps of $R$-modules.
Consider the pushout of $f$ and $g$,
$$
\xymatrix{
M  \ar[r]_f \ar[d]_g & N \ar[d]^{g'} \\
M' \ar[r]^{f'} & N'
}
$$
Then $g$ dominates $f$ if and only if $f'$ is universally injective.
\end{lemma}

\begin{proof}
Recall that $N'$ is $M' \oplus N$ modulo the submodule consisting of elements
$(g(x), -f(x))$ for $x \in M$.
From the construction of $N'$ we have a short exact sequence
$$
0 \to \Ker(f) \cap \Ker(g) \to \Ker(f) \to \Ker(f')
\to 0.
$$
Since tensoring commutes with taking pushouts, we have such a short exact
sequence
$$
0 \to \Ker(f \otimes \text{id}_Q ) \cap \Ker(g \otimes
\text{id}_Q) \to \Ker(f \otimes \text{id}_Q)
\to \Ker(f' \otimes \text{id}_Q) \to 0
$$
for every $R$-module $Q$.  So $f'$ is universally injective if and only if
$\Ker(f \otimes \text{id}_Q ) \subset \Ker(g \otimes
\text{id}_Q)$ for every $Q$, if and only if $g$ dominates $f$.
\end{proof}

\noindent
The above definition of domination is sometimes related to the usual notion
of domination of maps as the following lemma shows.

\begin{lemma}
\label{lemma-domination}
Let $f: M \to N$ and $g: M \to M'$ be maps of $R$-modules.
Suppose $\Coker(f)$ is of finite presentation.  Then $g$ dominates $f$ if
and
only if $g$ factors through $f$, i.e.\ there exists a module map $h: N
\to M'$ such that $g = h \circ f$.
\end{lemma}

\begin{proof}
Consider the pushout of $f$ and $g$ as in the statement of
Lemma \ref{lemma-domination-universally-injective}.
From the construction of the pushout it follows that
$\Coker(f') = \Coker(f)$, so $\Coker(f')$ is of finite
presentation.  Then by
Lemma \ref{lemma-universally-exact-split}, $f'$ is universally injective if and
only if
$$
0 \to M' \xrightarrow{f'} N' \to \Coker(f') \to 0
$$
splits. This is the case if and only if there is a map $h' : N' \to M'$
such that $h' \circ f' = \text{id}_{M'}$.  From the universal
property of the pushout, the existence of such an $h'$ is equivalent to $g$
factoring through $f$.
\end{proof}


\begin{proposition}
\label{proposition-ML-characterization}
Let $M$ be an $R$-module.  Let $(M_i, f_{ij})$ be a directed system of finitely
presented $R$-modules, indexed by $I$, such that $M = \colim M_i$.  Let
$f_i:
M_i \to M$ be the canonical map.  The following are equivalent:
\begin{enumerate}
\item For every finitely presented $R$-module $P$ and module map $f: P
\to M$, there exists a finitely presented $R$-module $Q$ and a module
map $g: P \to Q$ such that $g$ and $f$ dominate each other, i.e.,
$\Ker(f \otimes_R \text{id}_N) = \Ker(g \otimes_R \text{id}_N)$
for every $R$-module $N$.
\item For each $i \in I$, there exists $j \geq i$ such that $f_{ij}: M_i
\to M_j$ dominates $f_i: M_i \to M$.
\item For each $i \in I$, there exists $j \geq i$ such that $f_{ij}: M_i
\to M_j$ factors through $f_{ik}: M_i \to M_k$ for all $k \geq
i$.
\item For every $R$-module $N$, the inverse system
$(\Hom_R(M_i, N), \Hom_R(f_{ij}, N))$ is Mittag-Leffler.
\item For $N = \prod_{s \in I} M_s$, the inverse system
$(\Hom_R(M_i, N), \Hom_R(f_{ij}, N))$ is Mittag-Leffler.
\end{enumerate}
\end{proposition}

\begin{proof}
First we prove the equivalence of (1) and (2).  Suppose (1) holds and let $i
\in I$. Corresponding to the map $f_i: M_i \to M$, we can choose $g:
M_i \to Q$ as in (1).  Since $M_i$ and $Q$ are of finite presentation,
so is $\Coker(g)$.  Then by Lemma \ref{lemma-domination}, $f_i : M_i
\to M$ factors through $g: M_i \to Q$, say $f_i = h \circ g$
for some $h: Q \to M$.  Then since $Q$ is finitely presented, $h$
factors through $M_j \to M$ for some $j \geq i$, say $h = f_j \circ h'$
for some $h': Q \to M_j$.  In total we have a commutative diagram
$$
\xymatrix{
  &  M  & \\
M_i \ar[dr]_g \ar[ur]^{f_i} \ar[rr]^{f_{ij}} &
& M_j \ar[ul]_{f_j} \\
  & Q  \ar[ur]_{h'} &
}
$$
Thus $f_{ij}$ dominates $g$.  But $g$ dominates $f_i$, so $f_{ij}$ dominates
$f_i$.

\medskip\noindent
Conversely, suppose (2) holds.  Let $P$ be of finite presentation and $f: P
\to M$ a module map.  Then $f$ factors through $f_i: M_i \to M$
for some $i \in I$, say $f = f_i \circ g'$ for some $g': P \to M_i$.
Choose by (2) a $j \geq i$ such that $f_{ij}$ dominates $f_i$.  We have a
commutative diagram
$$
\xymatrix{
P \ar[d]_{g'} \ar[r]^{f}           & M  \\
M_i \ar[ur]^{f_i} \ar[r]_{f_{ij}} & M_j \ar[u]_{f_j}
}
$$
From the diagram and the fact that $f_{ij}$ dominates $f_i$, we find that $f$
and $f_{ij} \circ g'$ dominate each other.  Hence taking $g = f_{ij} \circ g' :
P \to M_j$ works.

\medskip\noindent
Next we prove (2) is equivalent to (3).  Let $i \in I$.  It is always true that
$f_i$ dominates $f_{ik}$ for $k \geq i$, since $f_i$ factors through
$f_{ik}$.  If (2) holds, choose $j \geq i$ such that $f_{ij}$ dominates
$f_i$.  Then since domination is a transitive relation, $f_{ij}$ dominates
$f_{ik}$ for $k \geq i$. All $M_i$ are of finite presentation, so
$\Coker(f_{ik})$ is of finite presentation for $k \geq i$.  By Lemma
\ref{lemma-domination}, $f_{ij}$ factors through $f_{ik}$ for all $k \geq i$.
Thus (2) implies (3).  On the other hand, if (3) holds then for any $R$-module
$N$, $f_{ij} \otimes_R \text{id}_N$ factors through $f_{ik}
\otimes_R \text{id}_N$ for $k \geq i$.  So $\Ker(f_{ik} \otimes_R
\text{id}_N) \subset \Ker(f_{ij} \otimes_R \text{id}_N)$ for $k
\geq i$.  But $\Ker(f_i \otimes_R \text{id}_N: M_i \otimes_R N
\to M \otimes_R N)$ is the union of $\Ker(f_{ik} \otimes_R
\text{id}_N)$ for $k \geq i$.  Thus $\Ker(f_i \otimes_R
\text{id}_N) \subset \Ker(f_{ij} \otimes_R \text{id}_N)$ for
any $R$-module $N$, which by definition means $f_{ij}$ dominates $f_i$.

\medskip\noindent
It is trivial that (3) implies (4) implies (5).  We show (5) implies (3).  Let
$N = \prod_{s \in I} M_s$. If (5) holds, then given $i \in I$ choose $j \geq i$
such that
$$
\Im( \Hom(M_j, N) \to  \Hom(M_i, N)) =
\Im( \Hom(M_k, N) \to  \Hom(M_i, N))
$$
for all $k \geq j$.  Passing the product over $s \in I$ outside of the
$\Hom$'s
and looking at the maps on each component of the product, this says
$$
\Im( \Hom(M_j, M_s) \to  \Hom(M_i, M_s)) =
\Im( \Hom(M_k, M_s) \to   \Hom(M_i, M_s))
$$
for all $k \geq j$ and $s \in I$.  Taking $s = j$ we have
$$
\Im( \Hom(M_j, M_j) \to  \Hom(M_i, M_j)) =
\Im( \Hom(M_k, M_j) \to   \Hom(M_i, M_j))
$$
for all $k \geq j$.  Since $f_{ij}$ is the image of
$\text{id} \in  \Hom(M_j, M_j)$ under
$\Hom(M_j, M_j) \to  \Hom(M_i, M_j)$,
this shows  that for any $k \geq j$ there is $h \in \Hom(M_k, M_j)$
such that $f_{ij} = h \circ f_{ik}$.  If $j \geq k$ then we can take
$h = f_{kj}$. Hence (3) holds.
\end{proof}

\begin{definition}
\label{definition-mittag-leffler-module}
Let $M$ be an $R$-module.  We say that $M$ is {\it Mittag-Leffler} if the
equivalent conditions of
Proposition \ref{proposition-ML-characterization}
hold.
\end{definition}

\noindent
In particular a finitely presented module is Mittag-Leffler.

\begin{remark}
\label{remark-flat-ML}
Let $M$ be a flat $R$-module. By Lazard's theorem
(Theorem \ref{theorem-lazard})
we can write $M = \colim M_i$ as the colimit of a
directed system $(M_i, f_{ij})$ where the $M_i$ are free
finite $R$-modules. For $M$ to be Mittag-Leffler, it is enough for the inverse
system of duals $(\Hom_R(M_i, R), \Hom_R(f_{ij}, R))$ to be
Mittag-Leffler. This follows from criterion (4) of
Proposition \ref{proposition-ML-characterization}
and the fact that for a free finite $R$-module $F$,
there is a functorial isomorphism
$\Hom_R(F, R) \otimes_R N \cong \Hom_R(F, N)$
for any $R$-module $N$.
\end{remark}

\begin{lemma}
\label{lemma-tensor-ML-modules}
If $R$ is a ring and $M$, $N$ are Mittag-Leffler modules over $R$,
then $M \otimes_R N$ is a Mittag-Leffler module.
\end{lemma}

\begin{proof}
Write $M = \colim_{i \in I} M_i$ and $N = \colim_{j \in J} N_j$
as directed colimits of finitely presented $R$-modules.
Denote $f_{ii'} : M_i \to M_{i'}$ and $g_{jj'} : N_j \to N_{j'}$ the
transition maps. Then $M_i \otimes_R N_j$ is a finitely presented
$R$-module (see
Lemma \ref{lemma-tensor-finiteness}),
and $M \otimes_R N = \colim_{(i, j) \in I \times J} M_i \otimes_R M_j$.
Pick $(i, j) \in I \times J$. By the definition of a Mittag-Leffler module
we have
Proposition \ref{proposition-ML-characterization} (3)
for both systems. In other words there exist $i' \geq i$ and $j' \geq j$
such that for every choice of $i'' \geq i$ and $j'' \geq j$ there exist
maps $a : M_{i''} \to M_{i'}$ and $b : M_{j''} \to M_{j'}$ such that
$f_{ii'} = a \circ f_{ii''}$ and $g_{jj'} = b \circ g_{jj''}$.
Then it is clear that
$a \otimes b : M_{i''} \otimes_R N_{j''} \to M_{i'} \otimes_R N_{j'}$
serves the same purpose for the system
$(M_i \otimes_R N_j, f_{ii'} \otimes g_{jj'})$.
Thus by the characterization
Proposition \ref{proposition-ML-characterization} (3)
we conclude that $M \otimes_R N$ is Mittag-Leffler.
\end{proof}

\begin{lemma}
\label{lemma-ML-also}
Let $R$ be a ring and $M$ an $R$-module. Then $M$ is Mittag-Leffler if and
only if for every finite free $R$-module $F$ and module map
$f: F \to M$, there exists a finitely presented $R$-module $Q$
and a module map $g : F \to Q$ such that $g$ and $f$ dominate each other, i.e.,
$\Ker(f \otimes_R \text{id}_N) = \Ker(g \otimes_R \text{id}_N)$
for every $R$-module $N$.
\end{lemma}

\begin{proof}
Since the condition is clear weaker than condition (1) of
Proposition \ref{proposition-ML-characterization}
we see that a Mittag-Leffler module satisfies the condition.
Conversely, suppose that $M$ satisfies the condition and that
$f : P \to M$ is an $R$-module map from a finitely presented
$R$-module $P$ into $M$. Choose a surjection $F \to P$ where
$F$ is a finite free $R$-module. By assumption we can find a map
$F \to Q$ where $Q$ is a finitely presented $R$-module such that
$F \to Q$ and $F \to M$ dominate each other. In particular, the kernel
of $F \to Q$ contains the kernel of $F \to P$, hence we obtain an
$R$-module map $g : P \to Q$ such that $F \to Q$ is equal to
the composition $F \to P \to Q$. Let $N$ be any $R$-module and
consider the commutative diagram
$$
\xymatrix{
F \otimes_R N \ar[d] \ar[r] & Q \otimes_R N \\
P \otimes_R N \ar[ru] \ar[r] & M \otimes_R N
}
$$
By assumption the kernels of $F \otimes_R N \to Q \otimes_R N$
and $F \otimes_R N \to M \otimes_R N$ are equal. Hence, as
$F \otimes_R N \to P \otimes_R N$ is surjective, also the kernels
of $P \otimes_R N \to Q \otimes_R N$
and $P \otimes_R N \to M \otimes_R N$ are equal.
\end{proof}

\begin{lemma}
\label{lemma-restrict-ML-modules}
Let $R \to S$ be a finite and finitely presented ring map.
Let $M$ be an $S$-module.
If $M$ is a Mittag-Leffler module over $S$ then
$M$ is a Mittag-Leffler module over $R$.
\end{lemma}

\begin{proof}
Assume $M$ is a Mittag-Leffler module over $S$.
Write $M = \colim M_i$ as a directed colimit of finitely presented
$S$-modules $M_i$. As $M$ is Mittag-Leffler over $S$ there exists for each
$i$ an index $j \geq i$ such that for all $k \geq j$ there is a factorization
$f_{ij} = h \circ f_{ik}$ (where $h$ depends on $i$, the choice of $j$ and
$k$). Note that by
Lemma \ref{lemma-finite-finitely-presented-extension}
the modules $M_i$ are also finitely presented as $R$-modules. Moreover, all
the maps $f_{ij}, f_{ik}, h$ are maps of $R$-modules. Thus we see that the
system $(M_i, f_{ij})$ satisfies the same condition when viewed as a system
of $R$-modules. Thus $M$ is Mittag-Leffler as an $R$-module.
\end{proof}

\begin{lemma}
\label{lemma-mod-ideal-ML-modules}
Let $R$ be a ring.
Let $S = R/I$ for some finitely generated ideal $I$.
Let $M$ be an $S$-module.
Then $M$ is a Mittag-Leffler module over $R$ if and only if
$M$ is a Mittag-Leffler module over $S$.
\end{lemma}

\begin{proof}
One implication follows from
Lemma \ref{lemma-restrict-ML-modules}.
To prove the other, assume $M$ is Mittag-Leffler as an $R$-module.
Write $M = \colim M_i$ as a directed colimit of finitely presented
$S$-modules. As $I$ is finitely generated, the ring $S$ is finite and finitely
presented as an $R$-algebra, hence the modules $M_i$ are finitely
presented as $R$-modules, see
Lemma \ref{lemma-finite-finitely-presented-extension}.
Next, let $N$ be any $S$-module. Note that for each $i$ we have
$\Hom_R(M_i, N) = \Hom_S(M_i, N)$ as $R \to S$ is surjective.
Hence the condition that the inverse system
$(\Hom_R(M_i, N))_i$ satisfies Mittag-Leffler, implies that the system
$(\Hom_S(M_i, N))_i$ satisfies Mittag-Leffler. Thus $M$ is
Mittag-Leffler over $S$ by definition.
\end{proof}

\begin{remark}
\label{remark-go-up-ML-modules}
Let $R \to S$ be a finite and finitely presented ring map.
Let $M$ be an $S$-module which is Mittag-Leffler as an $R$-module.
Then it is in general not the case that $M$ is Mittag-Leffler as
an $S$-module. For example suppose that $S$ is the ring of dual numbers
over $R$, i.e., $S = R \oplus R\epsilon$ with $\epsilon^2 = 0$. Then an
$S$-module consists of an $R$-module $M$ endowed with a square zero
$R$-linear endomorphism $\epsilon : M \to M$. Now suppose that $M_0$
is an $R$-module which is not Mittag-Leffler. Choose a presentation
$F_1 \xrightarrow{u} F_0 \to M_0 \to 0$ with $F_1$ and $F_0$ free $R$-modules.
Set $M = F_1 \oplus F_0$ with
$$
\epsilon =
\left(
\begin{matrix}
0 & 0 \\
u & 0
\end{matrix}
\right) : M \longrightarrow M.
$$
Then $M/\epsilon M \cong F_1 \oplus M_0$ is not Mittag-Leffler over
$R = S/\epsilon S$, hence not Mittag-Leffler over $S$ (see
Lemma \ref{lemma-mod-ideal-ML-modules}).
On the other hand, $M/\epsilon M = M \otimes_S S/\epsilon S$ which would
be Mittag-Leffler over $S$ if $M$ was, see
Lemma \ref{lemma-tensor-ML-modules}.
\end{remark}


